\begin{abstract}
Large-scale unsupervised language models (LMs) can acquire extensive world knowledge and certain reasoning capabilities, but precisely directing their outputs remains a challenge because their training lacks supervision. Current approaches to improve controllability typically involve gathering human judgements comparing model outputs and then adapting the pretrained LM to adhere to these preferences, frequently through reinforcement learning from human feedback (RLHF). Nonetheless, RLHF introduces significant complexity and instability, as it requires first training a reward model that captures human preferences, followed by reinforcement learning to tune the original LM towards maximizing this estimated reward, while constraining divergence from the pretrained model. In this work, we present a novel formulation of the reward model in RLHF that facilitates direct extraction of the optimal policy in closed form. This formulation allows the standard RLHF objective to be addressed via a straightforward classification loss. We introduce an algorithm called \textit{Direct Preference Optimization} (DPO), which offers greater stability, strong performance, and is computationally efficient; it does away with the need for sampling during fine-tuning or extensive hyperparameter searches. Through empirical evaluation, we demonstrate that DPO is capable of aligning LMs with human preferences on par with or surpassing contemporary techniques. In particular, DPO outperforms PPO-based RLHF methods for sentiment control, and provides equal or better results for response quality in summarization and single-turn dialogue tasks, all while reducing implementation and training complexity.
\end{abstract}

\section{Introduction}
Unsupervised language models (LMs) that are trained on massive corpora have demonstrated a range of unexpected abilities~\citep{chowdhery2022palm, brown2020language, touvron2023llama,bubeck2023sparks}. These models are exposed to data reflecting the diverse intentions, preferences, and proficiencies of human authors. Not all such attributes are worth emulating; for instance, it is advantageous for an AI coding assistant to \textit{recognize} common programming errors for the sake of correction, yet during code generation, we would prefer that the model emulate the exceptional—albeit less frequent—coding expertise available in its training examples. Similarly, while a model should \textit{know} about widespread misconceptions endorsed by 50\% of individuals, we certainly do not intend for it to present those misconceptions as facts in half of the relevant responses. Thus, discriminating between the model's \emph{preferred outputs and actions} and its broad \textit{scope of knowledge and skills} becomes fundamental to developing AI systems that are reliable, effective, and governable~\citep{ouyang2022training}. While current techniques primarily use reinforcement learning (RL) to align LMs with human preferences, we demonstrate that the conventional RL-based objective can, in fact, be equivalently optimized using a basic binary cross-entropy loss, streamlining the preference alignment workflow significantly.

Broadly speaking, current methodologies seek to instill desirable traits in language models by leveraging carefully selected human preference data that exemplifies qualities such as safety and helpfulness. This alignment process is executed subsequent to the initial phase of extensive unsupervised pre-training on textual data. Although the simplest strategy for incorporating preferences involves supervised fine-tuning with demonstrations of superior responses, reinforcement learning from human or AI feedback (RLHF/RLAIF; \citep{christiano2017deep,bai2022constitutional}) has emerged as the most effective paradigm. RLHF works by constructing a reward model based on annotated human preferences, followed by applying RL to update the language model policy to prioritize outputs with higher reward signals, all while constraining deviations from the pre-trained model. Despite the strong performance of RLHF in dialogue and code generation tasks, its implementation introduces considerable complexity beyond standard supervised learning—necessitating multiple rounds of language model training, as well as in-the-loop policy sampling, both of which entail substantial computational overhead.

This paper introduces a method for aligning language model behavior with human preferences by directly optimizing the model without the need for explicit reward modeling or reinforcement learning. We present \textit{{Direct Preference Optimization} (DPO)}, an approach that effectively targets the same optimization problem as RLHF methods—reward maximization constrained by KL-divergence—yet offers a more straightforward implementation and training process. The key intuition behind DPO is to enhance the log odds in favor of preferred responses over less preferred ones, but, crucially, it incorporates a dynamic importance weight for each example, thereby mitigating the model collapse issues encountered when using a simple probability ratio as the objective. DPO, similar to previous strategies, depends on a theoretical preference model (for example, the Bradley-Terry model; \cite{bradley1952rankanalysis}) to quantify the agreement between a reward function and observed preference data. In contrast to conventional pipelines—which employ the preference model to guide the training of a separate reward model, and subsequently optimize a policy with respect to that learned reward—DPO leverages a variable substitution to express the preference loss directly in terms of the policy itself. Provided a dataset that encodes human preference judgments on model outputs, DPO enables policy optimization with a standard binary cross-entropy loss, resulting in a policy that is optimal for an implicit reward function inferred from the preference data.

Our primary contribution is the proposal of Direct Preference Optimization (DPO), a reinforcement-learning-free, streamlined algorithm for training language models using human preferences. Empirically, we find that DPO matches or exceeds the performance of existing approaches, including PPO-based RLHF, in applications such as sentiment transfer, summarization, and conversational tasks with language models scaling up to 6B parameters.

\section{Related Work}

Progressively larger self-supervised language models have demonstrated the ability to solve certain tasks in a zero-shot setting \citep{radford2019language} or by leveraging few-shot prompting techniques \citep{gpt3,megatron,chowdhery2022palm}. Nevertheless, substantial gains in downstream performance and closer alignment with user objectives are obtained when these models are fine-tuned on collections of explicit instructions accompanied by human-authored completions \citep{mishra-etal-2022-cross,sanh2022multitask,chung2022scaling,thoppilan2022lamda}. This process, referred to as `instruction-tuning', enhances large language models’ ability to follow previously unseen instructions and improves their overall utility \citep{chung2022scaling}. Despite the demonstrated benefits of instruction tuning, it is frequently more practical to amass \textit{relative} assessments of output quality from humans than to acquire full expert demonstrations. Consequently, subsequent approaches have trained language models on human preference datasets, leading to better outcomes in tasks such as translation \citep{kreutzer-etal-2018-reliability}, summarization \citep{stiennon2022learning,ziegler2020finetuning}, story generation \citep{ziegler2020finetuning}, and instruction compliance \citep{ouyang2022training,ramamurthy2023is}. Typically, these techniques start by training a neural reward model that aligns with observed preferences, often operationalized using the Bradley-Terry model \citep{bradley1952rankanalysis}. Subsequently, the language model undergoes fine-tuning to maximize these rewards via reinforcement learning approaches such as REINFORCE \citep{williams1992reinforce}, proximal policy optimization (PPO; \cite{schulman2017proximal}), or related variants \citep{ramamurthy2023is}. In parallel, related research exploits instruction-tuned language models augmented with human feedback to produce synthetic datasets reflecting preferences over targeted criteria like safety and harmlessness \citep{bai2022constitutional}, employing human-supplied rubrics rather than detailed supervision to guide annotation. These efforts represent a merging of two established threads in the field: training language models via reinforcement learning for diverse goals~\citep{Ranzato2015SequenceLT,paulus2018a,wu2018learning}, and developing general strategies for learning from human preferences \citep{christiano2017deep,kupcsik2018learning}. Notwithstanding the conceptual attractiveness of using relative preferences, applying reinforcement learning to fine-tune large language models introduces considerable practical difficulties; in response, the present work puts forward a theoretically grounded alternative for optimizing relative preferences without recourse to reinforcement learning.

Beyond the domain of language, the problem of inferring policies from preference information has been addressed within both bandit and reinforcement learning frameworks, with multiple methodologies introduced. In contextual bandit scenarios where preferences or action rankings (instead of scalar rewards) guide learning, the problem is known as the contextual dueling bandit (CDB; \cite{yue2012karmed,dudik2015contextual}). In situations where absolute reward signals are missing, theoretical investigations into CDBs utilize the concept of a \textit{von Neumann winner}: a policy that achieves at least a 50\% expected win rate when compared to \textit{any} competing policy \citep{dudik2015contextual}. Nonetheless, CDBs are characterized by the online acquisition of preference feedback, whereas learning from human preferences in practice often involves a static, pre-collected set of preference-labeled action pairs \citep{yan2022human}. Likewise, in \textit{preference-based RL} (PbRL), the agent relies on binary preference signals derived from an \textit{unknown} underlying 'scoring' function, as opposed to direct rewards \citep{BusaFekete2014,ruiz2023dueling}. There exists a diverse array of PbRL algorithms, some of which leverage off-policy preference data; however, these approaches typically require an explicit estimation of the hidden scoring (reward) function before policy optimization commences \citep{jain2013learning,BusaFekete2014,christiano2017deep,sadigh2017active,kupcsik2018learning}. In contrast, we introduce a unified framework that learns a policy in a single stage by directly optimizing it to fulfill specified preferences.

\section{Preliminaries}\label{section:prelims}

We briefly outline the RLHF framework as presented in \citeauthor{ziegler2020finetuning} and extended by subsequent works \citep{stiennon2022learning, bai2022training, ouyang2022training}. This methodology typically encompasses three sequential components: (1) supervised fine-tuning (SFT); (2) preference elicitation and reward modeling; and (3) reinforcement learning-based optimization.

\textbf{SFT}: The standard RLHF workflow initiates with the supervised fine-tuning of a pre-trained language model using high-quality, task-specific datasets (such as for dialogue or summarization), yielding a model denoted as $\pi^\text{SFT}$.

\textbf{Reward Modeling Phase}: During this stage, the SFT model generates response pairs for prompts $x$, yielding $(y_1, y_2)\sim \pi^\text{SFT}(y \mid x)$. These answer pairs are then evaluated by human annotators, who state a preference for one completion over the other, represented as $y_w\succ y_l \mid x$, where $y_w$ is the selected (preferred) output and $y_l$ the less favored alternative. These preferences are believed to originate from an underlying, unobserved reward function $r^*(y, x)$. Several models have been devised to capture such preference structures, with the Bradley-Terry (BT) \cite{bradley1952rankanalysis} model being particularly prevalent (more broadly, the Plackett-Luce ranking model \citep{plackett1975analysis, luce2012individual} is applicable when multiple ordered completions are available). According to the BT model, the human preference probability $p^*$ takes the form:

\begin{equation}\label{eq:bradley-terry}
    p^*(y_1\succ y_2 \mid x)=\frac{\exp\left(r^*(x, y_1)\right)}{\exp\left(r^*(x, y_1)\right) + \exp\left(r^*(x, y_2)\right)}.
\end{equation}

Given a fixed dataset of pairwise comparisons $\mathcal{D}=\bigl\{x^{(i)}, y_w^{(i)}, y_l^{(i)}\bigr\}_{i=1}^N$ sampled from $p^*$, a reward function $r_{\phi}(x, y)$ can be introduced and its parameters fitted via maximum likelihood estimation. When this task is recast as binary classification, the corresponding negative log-likelihood objective becomes:

\begin{equation}\label{eq:reward_model}
    \mathcal{L}_R(r_{\phi}, \mathcal{D}) = -\mathbb{E}_{(x, y_w, y_l)\sim \mathcal{D}}\bigl[\log \sigma(r_{\phi}(x, y_w)- r_{\phi}(x, y_l))\bigr]
\end{equation}

where $\sigma$ is the sigmoid (logistic) activation function. For language models, the typical procedure is to initialize $r_{\phi}(x, y)$ from the SFT model $\pi^\text{SFT}(y \mid x)$, then augment the architecture with a linear projection atop the transformer’s output to yield a scalar reward prediction \cite{ziegler2020finetuning}. To promote stability in the reward estimation, normalization is commonly applied so that $\mathbb{E}_{x,y\sim \mathcal{D}}\left[r_\phi(x, y)\right] = 0$ for each $x$.

\textbf{RL Fine-Tuning Phase}: In the RL fine-tuning stage, the estimated reward model is used as a signal to adapt the language model policy. Consistent with prior studies~\citep{jaques2017sequence, jaques2020human}, the training objective can be expressed as:

\begin{equation}\label{eq:RL}
\max_{\pi_{\theta}}  \mathbb{E}_{x\sim \mathcal{D}, y\sim \pi_{\theta}(y \mid x)}\bigl[r_{\phi}(x, y)\bigr] - \beta\mathbb{D}_{\textrm{KL}}\bigl[\pi_{\theta}(y\mid x)\mid \mid \pi_\text{ref}(y\mid x)\bigr],
\end{equation}

where the hyperparameter $\beta$ regulates the extent to which the optimized policy $\pi_\theta$ diverges from the reference policy $\pi_\text{ref}$ (which is typically the original SFT model, $\pi^\text{SFT}$). Generally, $\pi_\theta$ is also initialized as $\pi^\text{SFT}$. The KL constraint is vital—it restricts the policy’s divergence to regions where the reward model is reliable and mitigates undesirable behaviors such as reduced diversity or concentration on a single output. Because generation in language models involves discrete sequences, this optimization problem is non-differentiable, necessitating the use of reinforcement learning methods. The established technique \citep{ziegler2020finetuning, stiennon2022learning, bai2022training, ouyang2022training} involves defining the reward for learning as ${r(x, y) = r_{\phi}(x, y) -\beta (\log \pi_{\theta}(y\mid x) - \log \pi_\text{ref}(y\mid x))}$ and maximizing it via PPO \cite{schulman2017proximal}. 

\section{Direct Preference Optimization}\label{sec:DPO}

Recognizing the difficulties that reinforcement learning presents for large-scale language model fine-tuning, our objective is to introduce a streamlined method for policy improvement that uses preference data directly. Traditional RLHF approaches train a reward function and subsequently apply RL to optimize the policy; in contrast, our framework exploits a particular reward model form that allows us to solve for its optimal policy exactly, sidestepping the need for an iterative RL loop. As elaborated in the following sections, the principal idea is to use a closed-form relationship between reward models and their optimal policies, transforming the training objective from one on rewards to one on policies. This change-of-variable methodology eliminates the requirement for an explicit, separate reward model, yet continues to optimize with respect to established models of human preferences, such as the Bradley-Terry framework. In summary

\textbf{Derivation of the DPO objective.} Our derivation begins from the conventional reinforcement learning objective defined in Eq.~\ref{eq:RL}, using a general reward function $r$. As established in previous studies~\citep{peters2007reinforcement, peng2019advantage, korbak2022reinforcement, go2023aligning}, it follows directly that the optimizer for the KL-regularized reward maximization problem presented in Eq.~\ref{eq:RL} is given by:

\begin{equation}\label{eq:op_policy}
    \pi_r(y\mid x) = \frac{1}{Z(x)}\pi_\text{ref}(y\mid x)\exp\left(\frac{1}{\beta}r(x, y)\right),
\end{equation}

where $Z(x) =\sum_{y}\pi_\text{ref}(y\mid x)\exp\left(\frac{1}{\beta}r(x, y)\right)$ denotes the normalization constant, or partition function (refer to Appendix \ref{app:derivation1} for a detailed derivation). Notably, estimating the partition function $Z(x)$ remains computationally demanding—even when adopting the MLE estimator $r_{\phi}$ in place of the true reward $r^*$~\citep{korbak2022reinforcement, go2023aligning}—which limits the practical utility of this parametrization. Nevertheless, Eq.~\ref{eq:op_policy} can be rearranged to write the reward in terms of the policy $\pi_r$, the reference $\pi_\text{ref}$, and the partition function $Z(x)$. Taking logarithms of both sides yields, after further manipulation:

\begin{equation}\label{eq:main_eq}
    r(x,y) =\beta \log \frac{\pi_r(y\mid x)}{\pi_\text{ref}(y\mid x)} + \beta \log Z(x).
\end{equation}

Applying this reparameterization for the ground-truth reward $r^*$ and associated optimal policy $\pi^*$, we note that the Bradley-Terry model only depends on reward differences across completions: specifically, ${p^*(y_1 \succ y_2 \mid x) = \sigma(r^*(x, y_1) - r^*(x, y_2))}$. Substituting Eq.~\ref{eq:main_eq} for $r^*(x, y)$ into the preference model in Eq.~\ref{eq:bradley-terry}, we see the partition function cancels out, allowing human preference likelihood to be recast exclusively in terms of $\pi^*$ and $\pi_\text{ref}$. In consequence, under the Bradley-Terry model, the optimal RLHF policy $\pi^*$ obeys:

\begin{equation}\label{eq:objective}
    p^*(y_1\succ y_2 \mid x)=\frac{1}{1 + \exp\left(\beta \log \frac{\pi^*(y_2\mid x)}{\pi_\text{ref}(y_2\mid x)} - \beta \log \frac{\pi^*(y_1\mid x)}{\pi_\text{ref}(y_1\mid x)}\right)}
\end{equation}

Complete derivations are available in Appendix~\ref{app:derivation2}. Although the discussion centers on the Bradley-Terry model, analogous results for Plackett-Luce preference structures~\citep{plackett1975analysis, luce2012individual} are provided in Appendix~\ref{app:plackett_luce_models}.

Having formulated the human preference likelihood directly via the optimal policy, we can now pose a maximum likelihood estimation objective for a parametrized policy $\pi_\theta$. In direct analogy to the reward modeling maximum likelihood approach (cf. Eq.~\ref{eq:reward_model}), the policy-based loss is:

\begin{equation}\label{eq:optimum_model}
    \mathcal{L}_\text{DPO}(\pi_{\theta}; \pi_\text{ref}) = -\mathbb{E}_{(x, y_w, y_l)\sim \mathcal{D}}\left[\log \sigma \left(\beta \log \frac{\pi_{\theta}(y_w\mid x)}{\pi_\text{ref}(y_w\mid x)} - \beta \log \frac{\pi_{\theta}(y_l\mid x)}{\pi_\text{ref}(y_l\mid x)}\right)\right].
\end{equation}

In this formulation, we learn an implicit reward through a distinct parameterization, where the optimal policy coincides with $\pi_\theta$. Additionally, because this methodology is tantamount to optimizing a reparametrized Bradley-Terry model, it inherits certain theoretical guarantees, such as consistency provided that the preference data follows appropriate distributional assumptions \cite{bong2022generalized}. We examine further theoretical aspects of DPO and its connections to other approaches in Section~\ref{sec:theory}.

\textbf{How does the DPO update function?} To gain insight into the inner workings of DPO, it is instructive to study the gradient of the loss term $\mathcal{L}_\text{DPO}$. The derivative with respect to the parameters $\theta$ is given by:

\begin{multline*}\label{eq:gradient}
    \nabla_\theta \mathcal{L}_\text{DPO}(\pi_\theta;\pi_\text{ref}) = \\ -\beta\mathbb{E}_{(x, y_w, y_l) \sim \mathcal{D}} \bigg[\underbrace{\sigma(\hat{r}_\theta(x, y_l) - \hat{r}_\theta (x, y_w))}_\text{larger value if reward model misorders examples}\bigg[\underbrace{\nabla_\theta\log \pi(y_w \mid x)}_\text{promote $y_w$ in likelihood} - \underbrace{\nabla_\theta\log\pi(y_l \mid x)}_\text{demote $y_l$ in likelihood}\bigg]\bigg],
\end{multline*}

where the reward signal $\hat{r}_\theta(x, y) = \beta \log \frac{\pi_\theta(y \mid x)}{\pi_\text{ref}(y \mid x)}$ is implicitly induced by the language model $\pi_\theta$ relative to the reference policy $\pi_\text{ref}$ (for further elaboration, see Section~\ref{sec:theory}). Conceptually, the gradient of $\mathcal{L}_\text{DPO}$ encourages increasing the probability of preferred completions $y_w$ and suppressing the probability of less desirable completions $y_l$. Crucially, the contribution of each sample is modulated by the extent to which the reward model $\hat{r}_\theta$ incorrectly favors the dispreferred completion, scaled by $\beta$, thereby reflecting both the misranking and the influence of the KL regularization. Empirically, this form of weighting proves essential: omitting it, as in a simplistic variant of the algorithm, may lead to problematic behaviors in the language model (see Appendix Table~\ref{tab:unlikelihood_generations}).

\textbf{DPO overview.} The standard workflow for DPO consists of: (1) For each prompt $x$, generate responses $y_1, y_2 \sim \pi_\text{ref}(\cdot \mid x)$, then use human preference judgments to assemble the offline preference dataset $\mathcal{D} = \{x^{(i)}, y_w^{(i)}, y_l^{(i)}\}_{i=1}^N$; (2) train the language model $\pi_\theta$ by minimizing the DPO loss $\mathcal{L}_\text{DPO}$ with respect to the chosen $\pi_\text{ref}$, the dataset $\mathcal{D}$, and the specified value of $\beta$.

Practically, it is advantageous to make use of publicly accessible preference datasets rather than conducting new sampling and preference elicitation. Because these datasets are typically constructed with samples from $\pi^\text{SFT}$, we set $\pi_\text{ref} = \pi^\text{SFT}$ when such a model is available. In situations where $\pi^\text{SFT}$ is missing, we initialize $\pi_\text{ref}$ by maximizing the likelihood over preferred outputs ${(x, y_w)}$, i.e., ${\pi_\text{ref} = \argmax_{\pi}\mathbb{E}_{x, y_w \sim \mathcal{D}}\left[\log \pi(y_w \mid x)\right]}$. This initialization strategy serves to reduce distribution mismatch between the unavailable true reference and the $\pi_\text{ref}$ employed by DPO. Implementation specifics and hyperparameter configurations are provided in Appendix~\ref{app:implementation}.

\section{Theoretical Analysis of DPO}
This section delivers a deeper theoretical interpretation of DPO, supports its foundation, and clarifies how DPO addresses shortcomings seen in RLHF actor-critic approaches, such as PPO~\cite{schulman2017proximal}.

\label{sec:theory}

\subsection{Your Language Model Is Secretly a Reward Model}
DPO eliminates the need for explicit reward modeling or a separate RL optimization loop by recasting both into a unified maximum likelihood training objective. Specifically, the objective in Eq. \ref{eq:main_eq} matches the Bradley-Terry formulation with the reward function $r^*(x, y) = \beta \log\frac{\pi^*_\theta(y \mid x)}{\pi_\text{ref}(y \mid x)}$. Optimization proceeds directly on the parametric policy $\pi_\theta$, paralleling the reward model fitting process from Eq. \ref{eq:reward_model} after a variable transformation. In this section, we formally develop the underlying theory for this reparameterization, demonstrate its flexibility in representing the full range of reward models, and establish conditions for precisely recovering the optimal policy. We start by introducing an equivalence relation for reward functions.

\begin{definition}
We define two reward functions $r(x, y)$ and $r'(x, y)$ as equivalent if and only if ${r(x, y)-r'(x, y) = f(x)}$ for some function $f$.
\end{definition}

It can be readily verified that this defines an equivalence relation, which divides the set of reward functions into disjoint equivalence classes. Based on this, we present two lemmas:

\begin{lemma}\label{lemma:same_prefrence} Within the Plackett-Luce family, and in particular under the Bradley-Terry model, any two reward functions belonging to the same equivalence class induce identical preference distributions.
\end{lemma}

\begin{lemma}\label{lemma:same_policy}
    Within the setting of constrained RL, two reward functions from the same equivalence class produce the same optimal policy.
\end{lemma}

Proofs for these lemmas are direct and are given in Appendix \ref{app:lemma1}. The first lemma highlights a familiar issue of under-specification in Plackett-Luce models, as documented by \cite{plackett1975analysis}. Because of this ambiguity, it is typically necessary to introduce additional identifiability conditions to guarantee meaningful MLE solutions from Eq. \ref{eq:reward_model} \cite{bong2022generalized}. The second lemma establishes that all reward functions from a given equivalence class yield the same optimal policy, so for our purposes, it suffices to recover any representative of the optimal reward class. The theorem below, proved in Appendix~\ref{app:thm1}, formalizes this statement:

\begin{theorem}\label{thm:main}
    Provided certain mild assumptions hold, every reward class consistent with the Plackett-Luce model (including Bradley-Terry as a special case) can be captured by the reparameterization ${r(x, y) = \beta \log \frac{\pi(y\mid x)}{\pi_\text{ref}(y\mid x)}}$ for some model $\pi(y\mid x)$ and a given reference distribution $\pi_\text{ref}(y \mid x)$.
\end{theorem}

\begin{sproof}
    Let $r(x, y)$ denote any reward function, with its corresponding optimal policy $\pi_r(y \mid x)$ defined as in Eq. \ref{eq:op_policy}. We will demonstrate that there exists a reward function in the equivalence class of $r$ expressible using the stated reparameterization. Consider the following transformation:
\begin{equation}
    f(r; \pi_\text{ref}, \beta)(x, y) = r(x, y) - \beta\log\sum_{y}\pi_\text{ref}(y\mid x)\exp\left(\frac{1}{\beta}r(x, y)\right)
\end{equation}
This operator $f$ normalizes the reward function by subtracting the log-partition function of $\pi_r$ (which depends only on $x$). Since the normalization term is a function of $x$ alone, the output remains within the equivalence class of $r(x, y)$. Finally, substituting the right-hand side of Eq.~\ref{eq:main_eq}, which applies for any $r$, yields $f(r; \pi_\text{ref}, \beta)(x, y) = \beta \log \frac{\pi_r(y\mid x)}{\pi_\text{ref}(y\mid x)}$. Therefore, the operator $f$ maps $r$ to a representative of its equivalence class with the required form, establishing that the proposed reparameterization is sufficiently general for our purposes.
\end{sproof}

An alternative interpretation of Theorem~\ref{thm:main} is that it determines, within each equivalence class of reward functions, the particular representative chosen by the DPO reparameterization. Specifically, it selects the reward function that fulfills:

\begin{equation}\label{eq:lag_p}
     \sum_{y}\underbrace{\pi_\text{ref}(y\mid x)\exp\left(\frac{1}{\beta}r(x, y)\right)}_{=\pi(y\mid x)\text{, using Thm.~\ref{thm:main} reparam.}} = 1,
\end{equation}

meaning that $\pi(y\mid x)$ defines a legitimate probability distribution, where all probabilities are non-negative and their total is unity. Inspecting Eq.~\ref{eq:op_policy} reveals that Eq.~\ref{eq:lag_p} expresses the partition function associated with the optimal policy derived from the reward function $r(x, y)$. The central idea behind the DPO algorithm is the introduction of constraints within the otherwise flexible Plackett-Luce (and, more specifically, the Bradley-Terry) preference model family. This procedure retains the set of reward models that can be represented, while also ensuring that the optimal policy described in Eq.~\ref{eq:op_policy} becomes analytically manageable for any given prompt $x$.

\subsection{Instability of Actor-Critic Algorithms} Our framework also provides insights into why standard actor-critic approaches, such as PPO, may encounter stability issues in RLHF scenarios. Adhering to the RLHF process and focusing on the RL fine-tuning phase described in Section \ref{section:prelims}, we can draw analogies to the control as inference formalism \cite{levine2018reinforcement} as it applies to the constrained RL problem set out in \ref{eq:RL}. Letting $\pi_{\theta}(y\mid x)$ be our parameterized policy, we seek to minimize $\mathbb{D}_{\text{KL}}[\pi_{\theta}(y|x) \mid \mid \pi^*(y\mid x)]$, where $\pi^*$ denotes the optimal policy generated from Eq.~\ref{eq:optimum_model} using reward function $r_{\phi}(y, x)$. With some algebraic manipulation, this results in the following objective function:

\begin{equation}\label{eq:AC}
    \max_{\pi_{\theta}}\mathbb{E}_{\pi_{\theta}(y\mid x)}\bigg[\underbrace{r_{\phi}(x, y) -\beta\log\sum_{y}\pi_\text{ref}(y\mid x)\exp\left(\frac{1}{\beta}r_{\phi}(x, y)\right)}_{f(r_{\phi}, \pi_\text{ref}, \beta)} - \underbrace{\beta\log\frac{\pi_{\theta}(y\mid x)}{\pi_\text{ref}(y\mid x)}}_{\text{KL}}\bigg]
\end{equation}

The objective function addressed here is identical to that optimized in previous studies 
\citep{ziegler2020finetuning, stiennon2022learning, bai2022training, ouyang2022training}, specifically when utilizing the DPO-equivalent reward for the reward function family $r_{\phi}$. Under these conditions, the normalization component within $f(r_{\phi}, \pi_\text{ref}, \beta)$ can be viewed as the soft value function corresponding to the reference policy $\pi_\text{ref}$. Though this normalization factor does not modify the ultimate optimum of the objective, omitting it can result in a policy gradient with high variance, thereby hindering stable training. One approach to address this normalization is to employ a learned value function; however, this may present its own optimization challenges. Previous work has tackled this issue by normalizing rewards based on a baseline derived from human completions, effectively serving as a single-sample Monte Carlo approximation of the normalization factor. By contrast, the DPO reparameterization produces a reward function that does not necessitate such baseline estimates.

\section{Experiments}
In this section, we present empirical analyses examining how effectively DPO can train policies directly from preference data. Our initial investigation considers a controlled text generation environment to answer: how does DPO’s efficiency in balancing reward maximization and minimizing KL-divergence from the reference policy compare to established preference optimization methods, such as PPO? We subsequently test DPO on larger-scale models and more complex RLHF challenges, such as summarization and conversational tasks. Our experiments show that DPO generally matches or exceeds the performance of competitive baselines like PPO-based RLHF and approaches that select the highest-scoring trajectory out of $N$ samples from a learned reward, all with minimal hyperparameter adjustment. Prior to discussing our findings, we outline the experimental methodology; supplementary implementation details can be found in Appendix~\ref{app:exp_details}.

\textbf{Tasks.} We investigate three distinct open-ended text generation tasks in our experimental framework. For each task, the learning algorithms derive a policy from a preference dataset denoted by $\mathcal{D}=\bigl\{x^{(i)}, y_w^{(i)}, y_l^{(i)}\bigr\}_{i=1}^N$. In the case of \textbf{controlled sentiment generation}, $x$ serves as the initial segment of a movie review extracted from the IMDb corpus \cite{maas-EtAl:2011:ACL-HLT2011}, and the objective for the policy is to produce $y$ that exhibits positive sentiment. To facilitate systematic evaluation, we synthetically create pairs of preferences by leveraging a pre-trained sentiment classifier, ensuring that $p(\text{positive}\mid x,y_w)>p(\text{positive}\mid x,y_l)$. The SFT baseline is established by further training GPT-2-large on the training subset of IMDb reviews until model convergence (implementation specifics can be found in App~\ref{app:sentiment_details}). For the \textbf{summarization} task, $x$ represents a Reddit forum post, with the policy expected to summarize the main ideas in $y$. Consistent with established methodologies, we utilize the Reddit TL;DR summarization corpus \citep{volske-etal-2017-tl} alongside human preference data provided by \citeauthor{stiennon2022learning}. The SFT model employed in this context is fine-tuned on human-composed summaries of forum threads\footnote{\url{https://huggingface.co/CarperAI/openai_summarize_tldr_sft}}, utilizing the TRLX framework \citep{leandro_von_werra_2023_7790115} to implement RLHF. The human feedback data, originally collected by \citeauthor{stiennon2022learning}, was generated based on outputs from an alternative but similarly fine-tuned SFT model. In the final task, \textbf{single-turn dialogue}, $x$ corresponds to a user prompt ranging from technical questions in astrophysics to personal advice inquiries. Here, the policy is tasked with composing an informative and engaging response $y$; this task employs the Anthropic Helpful and Harmless dialogue dataset \citep{bai2022training}, which comprises 170k conversational exchanges between human users and an automated assistant. Each dialogue concludes with a pair of model-generated responses and a label indicating which one was preferred by a human annotator. Since a pre-trained SFT model is not available in this scenario, we fine-tune a generic language model exclusively using the preferred responses to create the SFT reference.

\textbf{Evaluation.} We employ two complementary evaluation strategies in our experiments. For the controlled sentiment generation scenario, where the ground-truth reward function (a sentiment classifier) is available, we characterize each algorithm’s performance through the set of attainable trade-offs between reward and KL-divergence relative to a reference policy; this enables direct measurement against the constrained reward maximization goal. Conversely, in practical applications where the ground-truth reward is inaccessible, we adopt a comparative metric, measuring an algorithm's \textit{win rate} against a baseline using GPT-4 as an automated stand-in for human judges. This approach is used to assess summary quality in summarization tasks and the helpfulness of responses in single-turn dialogue tasks. Specifically, for summarization, the baseline is the reference summary from the test set, while for dialogue, it is the human-preferred test response. Recognizing that previous work indicates large language models can be effective automatic evaluators \citep{Chen2023ExploringTU}, we supplement our assessment with a dedicated human subject study to validate our reliance on GPT-4 (see Sec.~\ref{sec:human-judgments}). Our analysis demonstrates a high correspondence between GPT-4 and human evaluations, often matching or exceeding agreement between different human annotators.

\textbf{Methods.} Alongside DPO, we benchmark a suite of established methods for training language models in line with human preferences. For instance, we apply zero-shot prompting with \textbf{GPT-J} \citep{gpt-j} for summarization, and 2-shot prompting with \textbf{Pythia-2.8B} \citep{biderman2023pythia} for dialogue. We further assess models trained via supervised fine-tuning (\textbf{SFT}), as well as \textbf{Preferred-FT}, which fine-tunes the model on the selected completion $y_w$—chosen by the SFT model in sentiment and summarization tasks, or a general LM in dialogue tasks. An additional supervised-style baseline is \textbf{Unlikelihood}~\citep{welleck2019neural}, which trains the policy to boost the likelihood of $y_w$ while suppressing that of $y_l$, optionally modulated by a coefficient $\alpha\in[0,1]$ applied to the unlikelihood penalty. Our study also covers \textbf{PPO} \citep{schulman2017proximal}, using a reward model trained on preference data, as well as \textbf{PPO-GT}, which serves as an idealized variant relying directly on the ground-truth reward function within the sentiment setting. In the sentiment experiments, we employ two versions of PPO-GT: an out-of-the-box implementation \cite{leandro_von_werra_2023_7790115}, and a customized version incorporating reward normalization and refined hyperparameter tuning for enhanced results (these refinements are also applied to standard PPO with learned rewards). As a further reference point, we utilize a \textbf{Best of $N$} baseline, wherein $N$ outputs are generated per prompt from the SFT model (or Preferred-FT in dialogue), and the response with the highest reward according to a learned preference model is selected. While this strategy often yields high-performing responses, its practicality is limited by the necessity of sampling $N$ outputs for each test query, rendering it computationally expensive for moderate values of $N$.

\subsection{How well can DPO optimize the RLHF objective?}

In standard RLHF frameworks, the reward maximization objective is subject to a KL constraint, which serves to maximize the reward while limiting policy divergence from a reference policy. Consequently, a meaningful comparison between algorithms must consider not only the absolute reward attained, but also the corresponding KL divergence; maximizing reward at the expense of substantially increased KL is often undesirable. Figure~\ref{fig:frontier-tldr-main} illustrates the relationship between reward and KL for several methods in the sentiment classification task. For each algorithm, we perform several training experiments, varying the hyperparameter controlling policy conservativeness across different runs (target KL $\in\{3,6,9,12\}$ for PPO, $\beta \in \{0.05,0.1,1,5\}$, $\alpha\in\{0.05,0.1,0.5,1\}$ for unlikelihood, and different random seeds for preferred-FT), leading to a total of 22 experimental runs. At intervals of every 100 training steps, we evaluate policies using a held-out set of test prompts, calculating both the mean reward according to the true reward function and the mean sequence-level KL divergence\footnote{The sum of the KL-divergences across all timesteps.} from the reference policy, denoted as $\text{KL}\left(\pi\mid \mid \pi_\text{ref}\right)$. The results indicate that DPO establishes the most efficient reward-KL frontier by a significant margin, obtaining high reward levels while maintaining low KL divergence. This outcome is particularly noteworthy for several reasons. To begin with, both DPO and PPO are designed to optimize the identical objective, yet DPO demonstrates a marked increase in efficiency; DPO's tradeoff curve in reward/KL strictly surpasses that of PPO. Furthermore, DPO defines a superior frontier to PPO \emph{even when PPO is granted access to ground-truth rewards} (PPO-GT).

\subsection{Can DPO scale to real preference datasets?}
\label{sec:dpo-real-datasets}
We proceed to assess DPO’s fine-tuning capabilities using tasks in summarization and single-turn dialogue. In the summarization setting, standard automatic metrics such as ROUGE often fail to align well with human judgments~\citep{stiennon2022learning}. Previous research indicates that leveraging PPO to fine-tune language models on human preferences leads to more satisfactory summary outputs. For our evaluation, we sample generated summaries from various methods on the test split of the TL;DR summarization dataset, and compute each method’s mean win rate when pitted against the reference completions in this split. Sampling for all approaches is performed across temperatures ranging from 0.0 to 1.0, and the resulting win rates are depicted in Figure~\ref{fig:frontier-tldr-main} (right). For this comparison, DPO, PPO, and Preferred-FT are each fine-tuned starting from the same GPT-J SFT model\footnote{\url{https://huggingface.co/CarperAI/openai_summarize_tldr_sft}}. Our analysis shows that DPO attains an average win rate near 61\% at a sampling temperature of 0.0, surpassing PPO’s optimal win rate of approximately 57\% at the same temperature. In addition, DPO attains a superior maximum win rate compared to the best-of-$N$ baseline. It is worth mentioning that DPO’s $\beta$ hyperparameter was not extensively optimized in this study, so DPO’s performance may be underestimated here. DPO also demonstrates substantially greater resilience to varying temperature, as PPO’s win rate can drop to that of the initial GPT-J model at elevated temperatures. Notably, Preferred-FT shows little improvement over the baseline SFT model. We further provide a direct human comparison of DPO and PPO in Section~\ref{sec:human-judgments}: DPO outputs sampled at temperature 0.25 are chosen over PPO’s (sampled at temperature 0) in 58\% of evaluations.

For single-turn dialogue tasks, we conduct evaluations using the subset of the Anthropic HH dataset \citep{bai2022training} test split, focusing on instances involving a single human-assistant exchange. To assess the different methods, we use the preferred completions provided in the test set as reference outputs and compute win rates under GPT-4 evaluation. Due to the lack of an established SFT model for this scenario, we initialize experiments with a pre-trained Pythia-2.8B model, subsequently fine-tune a reference model using Preferred-FT to ensure that generated completions remain within the model’s output distribution, and ultimately perform DPO training. Our experiments also include a comparison with the strongest result among 128 completions generated via Preferred-FT (as the Best of $N$ baseline shows performance saturation at $N=128$ for this dataset; see Appendix Figure~\ref{fig:best-of-n}), as well as with a 2-shot prompted version of the base Pythia-2.8B. DPO matches or surpasses these baselines for the optimal temperature settings selected per method. 

Additionally, we examine the performance of an RLHF model—trained via PPO on the Anthropic HH dataset\footnote{\url{https://huggingface.co/reciprocate/ppo_hh_pythia-6B}} and provided by a recognized source\footnote{\url{https://github.com/CarperAI/trlx/tree/main/examples/hh}}—but fail to identify any prompt or sampling temperature that enables it to outperform the original Pythia-2.8B base model. Considering our findings from TL;DR and noting that both PPO and the Best of 128 approaches are optimized for the same reward objective, we use the Best of 128 metric as an approximate stand-in for PPO-level results. Collectively, DPO stands out as the sole computationally tractable method to yield improvements over the preferred completions in Anthropic HH, achieving performance on par with or exceeding that of the resource-intensive Best of 128 baseline. Moreover, Figure~\ref{fig:dialogue-main} demonstrates that DPO achieves its optimal outcomes in relatively few training steps.

\subsection{Generalization to a new input distribution}

To deepen our comparison of PPO and DPO under conditions of distributional shift, we assess the policies obtained from our Reddit TL;DR summarization study on an alternative data source: news articles found in the test partition of the CNN/DailyMail dataset \citep{nallapati-etal-2016-abstractive}. For consistency, we adopt the best-performing sampling temperatures determined on TL;DR (0 and 0.25). The resulting performance metrics appear in Table~\ref{tab:ood}. Evaluation relies on the GPT-4 win rate relative to gold-standard reference summaries in the datasets, using an identical GPT-4 (C) prompt as in Reddit TL;DR except for substituting “forum post” with “news article.” DPO maintains a clear advantage over PPO even on this novel distribution. These findings provide preliminary evidence that DPO policies demonstrate generalization capabilities on par with, or exceeding, those of PPO, despite DPO’s lack of access to the additional unlabeled Reddit TL;DR prompts incorporated into PPO training.

\subsection{Validating GPT-4 judgments with human judgments}
\label{sec:human-judgments}

We carry out a human evaluation to ascertain the consistency and dependability of GPT-4's judgments, utilizing summaries from the TL;DR experiment assessed with two distinct GPT-4 prompts. The \textbf{GPT-4 (S)} (simple) prompt asks solely which summary most effectively conveys the essential information in a post. In contrast, the \textbf{GPT-4 (C)} (concise) prompt also incorporates conciseness as a criterion; this prompt is included as we observed that, with the \textbf{GPT-4 (S)} prompt, GPT-4 favors summaries that are longer and more redundant than those preferred by human raters. Full prompt wordings are available in Appendix~\ref{app:prompts}. We conduct three pairwise evaluations representing the range of model quality—choosing the top-performing (DPO, temp. 0.25), lowest-performing (PPO, temp. 1.0), and an intermediate (SFT, temp. 0.25) method—each compared against the greedy PPO baseline (its optimal sampling temperature). Our analysis indicates that across both prompts, GPT-4’s ratings align with those of human judges to a degree comparable to inter-annotator human agreement, reinforcing GPT-4's suitability as a human surrogate (multiple human assessments are obtained only for the DPO and PPO-1 settings due to constraints on annotator availability). On the whole, the \textbf{GPT-4 (C)} prompt delivers win rates that better reflect human preferences; as a result, we use this prompt in reporting main outcomes in Section~\ref{sec:dpo-real-datasets}. Further information on the human evaluation protocol, including the rater interface and the roster of human participants, is provided in Appendix~\ref{app:human-study}.

\section{Discussion}
Preference-based learning provides an effective and scalable approach for developing proficient and aligned language models. In this work, we presented DPO, a streamlined training method that utilizes preferences to guide language model development, circumventing the need for reinforcement learning. Unlike traditional approaches that adapt the preference learning challenge into an RL context merely to leverage standard RL algorithms, DPO establishes a correspondence between the policies of language models and reward functions. This correspondence facilitates direct alignment of models with human preferences by applying a straightforward cross-entropy loss, thereby avoiding reinforcement learning procedures without any reduction in expressive capacity. DPO matches or exceeds the performance of prevailing RLHF algorithms, such as those employing PPO, while requiring almost no hyperparameter adjustment, thereby lowering practical hurdles to preference-based language model training.

\textbf{Limitations \& Future Work.} Our findings prompt several avenues for further exploration. For instance, what are the generalization characteristics of DPO-trained policies in comparison to those optimized via explicit reward functions? Preliminary experiments indicate that DPO achieves generalization abilities akin to PPO-based methods, but further, more systematic investigation is warranted. Questions remain as to whether employing DPO self-labeling can take advantage of unlabeled prompts effectively. Additionally, the manifestation of reward over-optimization within the DPO paradigm merits investigation; specifically, the modest performance drop observed in Figure~\ref{fig:dialogue-main}-right could potentially reflect this phenomenon. While our evaluations currently extend to models with up to 6B parameters, scaling DPO to models of substantially greater size represents an important future research direction. In terms of assessment methodologies, we have observed that the prompt significantly influences the win rates measured by GPT-4, suggesting future work should focus on methods to obtain more reliable automated evaluations. Finally, DPO’s potential applications extend beyond preference alignment in language models and could include training of generative models across different modalities.